% =====================================================================
% preambulo.tex -- pacotes e comandos comuns. NAO precisa editar.
% =====================================================================
\usepackage[utf8]{inputenc}
\usepackage[T1]{fontenc}
\usepackage[brazil]{babel}
\usepackage{lmodern}
\usepackage[margin=2.3cm]{geometry}
\usepackage{amsmath,amssymb}
\usepackage{tikz}
\usetikzlibrary{arrows.meta,positioning}
\usepackage{booktabs,tabularx,array}
\usepackage{xcolor}
\usepackage{listings}
\usepackage{float}
\usepackage{caption}
\usepackage[hidelinks]{hyperref}
\setlength{\emergencystretch}{3em}

% ---------- cores dos blocos ----------
\definecolor{cora}{RGB}{140,90,255}
\definecolor{corb}{RGB}{255,240,0}
\definecolor{corc}{RGB}{0,170,255}
\definecolor{cord}{RGB}{235,30,30}

% ---------- blocos de codigo / saida de terminal ----------
\lstset{basicstyle=\ttfamily\footnotesize,frame=single,breaklines=true,
  columns=fullflexible,keepspaces=true,backgroundcolor=\color{black!3},
  rulecolor=\color{black!30}}

% ---------- marcadores para o trabalho em equipe ----------
% \afazer{texto}      -> aviso vermelho no PDF do que ainda falta escrever
% \responsavel{Nome}  -> quem escreve a secao (aparece em cinza)
% Para sumir com os dois na versao final, troque \marcadorestrue por
% \marcadoresfalse no main.tex.
\newif\ifmarcadores
\newcommand{\afazer}[1]{\ifmarcadores\par\noindent
  \textcolor{red!75!black}{\textbf{[A FAZER]} #1}\par\fi}
\newcommand{\responsavel}[1]{\ifmarcadores\par\noindent
  \textcolor{black!50}{\small\itshape Responsável: #1}\par\smallskip\fi}

% ---------- desenho de estados (TikZ) ----------
% Dentro de um tikzpicture:
%   \mesa                      desenha a mesa com os pontos 0..6
%   \bl{p}{l}{comprimento}{nome}   bloco comecando no ponto p, nivel l
%   \rot{texto}                rotulo acima do estado (ate o nivel 1)
%   \rotA{texto}               rotulo mais alto (estados com nivel 2)

% ---------- desenho de estados ----------
\newcommand{\mesa}{%
	\draw[line width=0.7pt] (0,0) -- (3,0);
	\foreach \i in {0,...,6}{\draw (\i*0.5,0) -- (\i*0.5,-0.07);
		\node[below,inner sep=1pt,font=\tiny] at (\i*0.5,-0.07) {\i};}}
% \bl{ponto p}{nivel l}{comprimento}{nome}
\newcommand{\bl}[4]{%
	\filldraw[fill=cor#4,draw=black!70] (#1*0.5,#2*0.5) rectangle ++(#3*0.5,0.5);
	\node[font=\small] at (#1*0.5+#3*0.25,#2*0.5+0.25) {\textit{#4}};}
\newcommand{\rot}[1]{\node[font=\small] at (1.5,1.35) {#1};}

\newcommand{\Szero}{\bl{0}{0}{2}{c}\bl{3}{0}{1}{a}\bl{5}{0}{1}{b}\bl{3}{1}{3}{d}}
\newcommand{\Sum}{\bl{0}{0}{2}{c}\bl{3}{0}{1}{a}\bl{5}{0}{1}{b}\bl{0}{1}{3}{d}}
\newcommand{\Sdois}{\bl{0}{0}{2}{c}\bl{0}{1}{3}{d}\bl{5}{0}{1}{b}\bl{5}{1}{1}{a}}
\newcommand{\Stres}{\bl{0}{0}{2}{c}\bl{2}{0}{3}{d}\bl{5}{0}{1}{b}\bl{5}{1}{1}{a}}
\newcommand{\Sfq}{\bl{0}{0}{2}{c}\bl{2}{0}{3}{d}\bl{5}{0}{1}{b}\bl{0}{1}{1}{a}}
\newcommand{\rotA}[1]{\node[font=\small] at (1.5,1.85) {#1};}
% Situacao 2
\newcommand{\Tzero}{\bl{0}{0}{2}{c}\bl{3}{0}{3}{d}\bl{0}{1}{1}{a}\bl{1}{1}{1}{b}}
\newcommand{\Tum}{\bl{0}{0}{2}{c}\bl{2}{0}{1}{b}\bl{3}{0}{3}{d}\bl{0}{1}{1}{a}}
\newcommand{\Tdois}{\bl{0}{0}{2}{c}\bl{2}{0}{1}{b}\bl{3}{0}{3}{d}\bl{2}{1}{1}{a}}
\newcommand{\Ttres}{\bl{2}{0}{1}{b}\bl{3}{0}{3}{d}\bl{2}{1}{1}{a}\bl{4}{1}{2}{c}}
\newcommand{\Tquatro}{\bl{2}{0}{1}{b}\bl{3}{0}{3}{d}\bl{4}{1}{2}{c}\bl{4}{2}{1}{a}}
\newcommand{\Tcinco}{\bl{3}{0}{3}{d}\bl{4}{1}{2}{c}\bl{4}{2}{1}{a}\bl{5}{2}{1}{b}}
% Situacao 3
\newcommand{\Ucinco}{\bl{0}{0}{2}{c}\bl{2}{0}{3}{d}\bl{0}{1}{1}{a}\bl{1}{1}{1}{b}}
\newcommand{\Usete}{\bl{0}{0}{2}{c}\bl{3}{0}{3}{d}\bl{0}{1}{1}{a}\bl{1}{1}{1}{b}}

% ---------- predicados (modo matematico) ----------
\newcommand{\at}{\mathit{at}}
\newcommand{\lev}{\mathit{lev}}
\newcommand{\clr}{\mathit{clr}}
\newcommand{\cov}{\mathit{cov}}
\newcommand{\livre}{\mathit{livre}}
\newcommand{\esp}{\mathit{esp}}
\newcommand{\on}{\mathit{on}}
\newcommand{\mv}{\mathit{move}}
\newcommand{\spn}{\mathit{span}}
\newcommand{\ov}{\mathit{overlap}}
\newcommand{\mesaT}{\mathsf{T}}
